\documentclass[10pt]{article}
\usepackage[margin=20mm]{geometry}
\usepackage{amsmath,amssymb,fontspec,unicode-math,xurl,hyperref}
\IfFontExistsTF{DejaVu Serif}{\setmainfont{DejaVu Serif}}{\setmainfont{Latin Modern Roman}}
\setmathfont{latinmodern-math.otf}
\setlength{\parindent}{0pt}\setlength{\parskip}{6pt}\sloppy
\begin{document}
Abstract: This initial review checks the proceedings paper, all embedded appendices, and a pinned released source snapshot. Two problems affect central links in the proposed method: narrowing a policy support does not imply the asserted entropy reduction, and the released tree selector omits ancestor factors required by the stated selection rule. The latter changes the chosen branch on an actual supplied tree. A separate local implementation defect lets evaluation mutate the training environment during PPO rollouts. The appendix’s discrete entropy support bound and its finite-policy identification argument withstand the checks described here. The empirical results support useful strategy-conditioned behavior within their tested settings, but historical code provenance, total interaction costs, and some selection details remain unresolved. No reinforcement-learning experiment or author program was executed.\par
Date: 24 October 2026\par
Title: Scientific review of Strategic Planning:
specificity, strategy selection, and rollout integrity\par
\# Identity, version, and review limits\par
The reviewed work is Max Ruiz Luyten, Antonin Berthon, and Mihaela van der Schaar, Strategic Planning: A Top-Down Approach to Option Generation, ICML 2025, PMLR 267, pp. 52258–52302, local number 0198, paper identifier I25-43567, and OpenReview identifier `xkgQWEj9F2` [Retained latex: \textbackslash{}cite\{paper\}]. Page references are one-based pages of the 45-page proceedings PDF, including its appendices. Its SHA-256 is\par
13ab03d69086576727bee5ef565d9970
41f5339fb2de47634ad0efefc9b2ccb2.\par
The 57-file released repository snapshot was inspected at commit |0321ac2cfb1b6a3b407299fc8ce41538b7f24004| [Retained latex: \textbackslash{}cite\{code\}]. The acquisition record does not establish this as the commit that generated the published experiments.\par
The review covers the main paper, Appendix A’s proofs, implementation and prompts in Appendix B, experimental tables and figure descriptions in Appendix C, related-work and illustrative discussions in Appendix D, and limitations in Appendix E. Relevant released code was read statically: tree construction and selection, prompts, reward-score collection and distillation, PPO reward mixing and rollout collection, evaluation callbacks, configurations, and supplied tree JSON. The existing PDF was actually read through Poppler; its path was not treated as evidence of access.\par
Separate OpenReview attachments remain unavailable in the acquisition record (HTTP 403). The embedded appendices were available and read. Raw seed-level learning curves, model checkpoints, annotation datasets, and the experimental EDE/Dreamer implementations are not present in the supplied source snapshot. Submodule contents were not downloaded. These limits prevent historical numerical reproduction and causal attribution of the published gains. They do not prevent the finite mathematical and static implementation checks below.\par
\# Contribution, definitions, and checked positive results\par
The paper proposes choosing a distribution \(\delta_\Pi\) over policies before low-level reinforcement learning. Its Strategy Problem trades off specificity, represented by entropy, and retention of a near-optimal policy. The implementation builds a natural-language strategy tree using an LLM, estimates each node’s value and feasibility, and converts a chosen strategy into a distilled state score used for reward shaping. Experiments compare strategy-guided and baseline PPO/EDE agents in original and modified Crafter.\par
For a discrete prior with policy masses \(p_i\), write \[H(p)=-\sum_i p_i\log p_i,\qquad
 S_\tau=\{\pi_i:p_i\geq\tau\}.\] The appendix explicitly restricts its entropy argument to finite or countable policy collections. It assumes that a near-optimal policy belongs to \(S_\tau\), that candidate policies can be executed, and that episodes are independently sampled from the same initial-state distribution. These are substantive requirements. A natural-language node and a neural policy learner do not automatically provide an enumerable thresholded policy set or the required distribution.\par
\#\# Entropy support bound\par
Equation reference: eq:support. Lemma A.1, pp. 14–15, correctly gives \[
 |S_\tau|\leq
 \min\left\{\frac{1}{\tau},
       2+\frac{H_{\max}}{\tau\log(1/\tau)}\right\}.\] The normalization bound follows from \(\tau|S_\tau|\leq1\). For the entropy bound, there are at most two atoms with mass greater than \(1/e\). On the remaining atoms, \(-x\log x\) is increasing up to \(1/e\), so each atom in \(S_\tau\) contributes at least \(\tau\log(1/\tau)\) when that light-atom set is nonempty. The constant two is therefore essential in the stated general bound and is present in the actual appendix. At \(\tau=1\), the appendix supplies the appropriate endpoint convention and at most one policy remains. A criticism that the proof incorrectly assumes all masses are at most \(1/e\) would be unfounded.\par
\#\# Policy identification\par
For bounded independent episodic returns in \([0,1]\), Proposition A.2 uses \[m=\left\lceil\frac{2}{\varepsilon^2}
                   \log\frac{4P}{\delta}\right\rceil\] episodes per policy. Hoeffding’s inequality gives simultaneous mean errors at most \(\varepsilon/2\) with failure probability at most \(\delta/2\), which is stronger than the claimed \(\delta\). Comparing the empirical maximizer with the true maximizer then gives an \(\varepsilon\) value gap.\par
For discounted rewards \(r_t\in[0,1]\), Lemma A.3 explicitly includes the return range \((1-\gamma)^{-1}\) and the truncation tail \(\gamma^T/(1-\gamma)\). Its chosen horizon makes the tail at most \(\varepsilon/2\), while its sampling rule bounds statistical error by another \(\varepsilon/2\). Theorem A.4 then uses value estimates accurate to \(v_{\min}/4\) to obtain algorithmic error \(v_{\min}/2\); adding the prior’s \(v_{\min}\) approximation gap gives the stated \(3v_{\min}/2\) result. This dependency and the discounted \((1-\gamma)^{-2}\) factor are consistent.\par
The compact main-text Theorem 2.1 must be read with these appendix restrictions, not as a return-scale-free guarantee for arbitrary unnormalized MDPs or a theorem about the implemented neural agent. The paper itself explicitly states on p. 4 that the theoretical bound does not extend to its implementation. No failure of Appendix A is asserted here.\par
\#\# Empirical and reward checks\par
Table 1 agrees with Table 5 for the short-schedule Crafter-Medium milestones. Table 2 agrees with Table 8 for long-schedule Crafter-Easy strategy metrics. The reported Crafter-Medium products in Table 9 are \(0.75(0.85)=0.6375\), \(0.9(0.5)=0.45\), and \(0.6(0.7)=0.42\), consistent with the printed ranks after rounding.\par
The data support a task- and strategy-dependent account, rather than uniform superiority at every milestone. In Tables 3–4, plain PPO reaches the easy hunting threshold at \(0.20\) million steps, before either guided strategy under either schedule. The medium results and the breeding behavior are stronger evidence for the intended difficult-exploration use case. This observation bounds broad prose about consistently faster learning; it does not negate the reported medium gains.\par
The released direct PPO reward mode computes a convex mixture of environment reward and a scaled distilled state score, consistent with Appendix B.4’s current-state formula and \(\beta=10\) through the configuration’s \(0.1\) scale. Such a score is not generally potential-based reward shaping, but the paper does not claim policy invariance during the shaping phase. Its explicit annealing and discussion of strategy reversion acknowledge that the optimization objective changes. The distinction between the next-state formula in the main text and current-state implementation should be specified in a reproduction, but no unsupported invariance violation is alleged.\par
\# Problem 1: a smaller policy support need not have lower entropy\par
\#\# Claim and suspect implication\par
Section 4.1, p. 5, constrains children to increasingly specific policy subsets. Section 4.3 then asserts \[H(\delta_{h_{i,j}})<H(\delta_{h_i})\] as a consequence of the tree’s design. Appendix B.2.1, p. 19, repeats that narrowing subgoals implicitly drives entropy down and explains that neither entropy nor the threshold \(H_{\max}\) is numerically checked. This is the proposed connection between the implementable tree construction and the formal specificity objective.\par
Support inclusion alone does not imply that Shannon entropy decreases. This remains false under the natural, favorable interpretation that a child distribution is exactly its parent’s conditional distribution on the retained policies, rather than an arbitrarily reweighted new distribution.\par
\#\# Complete counterexample\par
Consider a one-step finite MDP. Its initial state has three actions, each deterministically gives reward one and terminates. Fix \(\gamma=0.9\); its value is immaterial after termination. Let \(\pi_0,\pi_1,\pi_2\) be the deterministic policies choosing the respective actions. All policies are implementable, and each has return \(J(\pi_i)=1=J(\pi^\star)\).\par
Give the parent strategy masses \[p=(0.98,0.01,0.01).\] Let a more restrictive child require choosing either action one or action two. Its support is the proper subset \(\{\pi_1,\pi_2\}\). Condition the parent on that subset: \[q(\pi_1)=q(\pi_2)=\frac{0.01}{0.02}=\frac12.\] Both distributions are normalized, have finite Shannon entropy, and refer to policies in the same MDP. The refinement only removes a permitted policy. It preserves every remaining policy’s optimal value and implementability. Thus no loss of value, infeasibility, unknown dynamics, continuous density, or infinite support explains the example.\par
With natural logarithms, \[\begin{aligned}
 H(p)&=-0.98\log0.98-2(0.01\log0.01)
       =0.1119020569\ldots,\\
 H(q)&=-2\left(\tfrac12\log\tfrac12\right)
       =\log2=0.6931471806\ldots.\end{aligned}\] The support shrinks from three policies to two, but entropy increases by approximately \(0.5812451\) nats. Choose \(H_{\max}=0.2\), \(\tau=0.005\), and \(v_{\min}=0.1\). The parent satisfies the entropy budget; the child violates it. Every retained policy lies above the probability threshold and has value one, so both the relative and additive near-optimality requirements hold. Feasibility is one for both strategies, and the child-to-parent optimal-value ratio is one. All the operative value and feasibility conditions therefore survive while the claimed specificity implication fails.\par
The independent arithmetic check reproduces these numbers in |independent\_checks\_output.json|. More generally, \(H(p)=h(0.02)+0.02\log2\) here, whereas entropy conditional on the rare two-policy subset is \(\log2\). Conditioning can increase uncertainty for a particular selected branch, even though the average conditional entropy obeys a different inequality.\par
\#\# Alternative readings and repair scope\par
If every node were explicitly uniform on its finite support, strict support reduction would imply \(H=\log|\operatorname{supp}|\) decreases. An explicit entropy-contraction constraint would also suffice by assumption. Neither condition is imposed on the paper’s general strategy distributions or implemented by its language prompts. A verbal goal that is more concrete specifies no probability weights.\par
The assertion may be read as an intended design objective rather than a proven property. Under that reading, the formal inequality should be replaced by a heuristic hypothesis and should not be presented as guaranteed by hierarchy. Appendix B.2.1’s acknowledgment that entropy is not computed is important counterevidence to any claim that the implementation certifies the Strategy Problem; it does not make the displayed implication mathematically true.\par
A precise repair is to state and enforce an additional distributional condition, or to withdraw the entropy-contraction implication and describe specificity operationally. The uniform-support repair restores this particular implication analytically, but does not prove that natural-language refinements preserve near-optimal policies or accelerate PPO. No implementation of such a repair was tested.\par
\#\# Dependency and severity\par
The issue is Major. Specificity is one of the two defining quantities of the Strategy Problem, and automatic entropy reduction is the stated reason that the strategy tree addresses it without calculating entropy. The counterexample invalidates that central derivation. It does not invalidate the conditional PAC theorem or demonstrate that the recorded Crafter improvements disappear. Those surviving results do not remove the direct dependency of the formal-to-algorithm connection on the false implication.\par
\# Problem 2: the released selector omits ancestor factors\par
\#\# Stated rule and actual dataflow\par
Equation reference: eq:priority. Appendix B.2, p. 18, selects a leaf using its root-relative feasibility and value. For a path \((h_0,h_1,\ldots,h_m)\), the displayed formulas give \[
 Q(h_m)=\prod_{j=1}^m
         f(h_j\mid h_{j-1})\,V(h_j\mid h_{j-1}).\] This discounts a locally promising subgoal when its ancestors retain little feasibility or value.\par
In the pinned `strategist/tree.py`, lines 57–83, |find\_most\_promising\_leaf| performs a depth-first search and computes only\par
{\ttfamily\footnotesize return node["feasibility"] * node["value"]}\par
at each leaf or sequential node. The recursive return passes this score upward unchanged. No ancestor product is accumulated. Both strategy-tree refinement and default final selection call this method in `strategist\_agent.py`. The prompts define value relative to the parent, so the omitted factors cannot generally be absorbed into the stored value.\par
\#\# Check on an actual released tree\par
The supplied provides a decisive non-synthetic example. Treat sequential nodes as terminal candidates, exactly as the released selector does. Root feasibility and value both equal one. Each relevant field is a finite number in \([0,1]\). Applying the two explicitly specified arithmetic rules to the same JSON gives:\par
Candidate | Leaf-only score | Root-path score\par
Hunt, node 2 | \(0.4200\) | \(0.4200\)\par
Craft plan, node 6 | \(0.4650\) | \(0.2976\)\par
Breed plan, node 15 | \(0.4900\) | \(0.2646\)\par
Survival plan, node 11 | \(0.5200\) | \(0.3276\)\par
For example, the survival plan has local \(f=0.6\) and \(V=0.866666\ldots\), giving \(0.52\). Its parent has \(f=0.9\) and \(V=0.7\), so the paper’s score is \(0.9(0.7)(0.52)=0.3276\). Direct hunting is a root child with score \(0.7(0.6)=0.42\).\par
The released selection arithmetic therefore chooses survival (node 11); the stated root-relative rule chooses hunting (node 2). Both maxima are strict, so neither tie handling nor floating-point rounding explains the difference. The independent script reads only JSON data and computes both formulas; it does not import or execute `Tree`.\par
\#\# Counterevidence and historical scope\par
This mismatch does not change every tree’s best candidate. The corresponding reflected easy tree and both supplied medium variants retain the same top candidate under the two rules. Those controls are included in the evidence output. Further, command-line `–branch\_i` or `–node` choices can bypass automatic selection during training. The source’s `Config.tree\_file` points to `tree.json`; separate reflection files exist, but their role in each historical run is not reconstructed here.\par
If the stored fields were already root-relative, a leaf-only product could be appropriate. That interpretation conflicts with the parent-relative value description and with Equation [eq:priority] (\#eq:priority); it would require a different representation and an explicit conversion. It is not a resolution of the current paper–code discrepancy.\par
The paper also leaves a related selection provenance question. Section 5.2 says that the top two modified-Crafter strategies are selected by feasibility and value, but Tables 9–10 rank Craft, Hunt, and Breed in that order, while the main medium comparisons emphasize Hunt and Breed. Figure 3 and the post-generation reflection step provide possible evidence of different ranking stages. Without stage-specific run records, the discrepancy is recorded as Needs Evidence about the historical selection procedure, not as proof of post hoc cherry-picking.\par
\#\# Dependency, severity, and repair\par
The issue is Major for the released automatic selection implementation. The ranking rule directly determines which strategy branch is refined or chosen for reward shaping, a central algorithmic decision. The actual supplied tree demonstrates a different chosen strategy, not merely a numerically different unused field. The paper’s intended algorithm and its release therefore disagree at a core step.\par
Multiplying the local factors along each root-to-candidate path restores Equation [eq:priority] (\#eq:priority) for a fixed tree and declared parent-relative attributes; the calculation above verifies that specific repair. A log-domain implementation could avoid underflow. This repair does not validate LLM estimates, establish the historical experiment commit, or predict revised learning curves. Those would require separate provenance and experiments.\par
\# Problem 3: evaluation changes the training environment mid-rollout\par
\#\# Source claim and implementation\par
The experimental procedure and Appendix B.5 describe ordinary PPO training with periodic evaluation. In |strategist/rl\_agent.py|, `RLAgent.train`, lines 229–235, sets\par
{\ttfamily\footnotesize eval\_env = self.env}\par
and supplies that object to the callback. At each logging interval, |WandbCallback.\_on\_step|, lines 93–137, resets this environment and steps it through the evaluation episodes. The provided PPO configurations use a 20,000-step logging interval and ten evaluation episodes.\par
In `custom\_ppo.py`, the training environment step occurs at line 188. The callback runs at lines 193–195. Only afterward does the collector store the transition and assign the saved `new\_obs` to `self.\_last\_obs`, at lines 220–229. There is no restoration of the environment’s state or synchronization of the learner’s cached observation after evaluation. Passing the same object creates an actual state mutation, even if a vector wrapper surrounds it for training.\par
\#\# Complete mechanism check\par
Let \(s_{t+1}\) be the state just reached by a training step. The collector saves its observation. The callback then resets the same underlying environment and completes evaluation trajectories, leaving its state at \(s_{\rm eval}\). The next training action is chosen using the saved observation of \(s_{t+1}\) but applied to \(s_{\rm eval}\).\par
An independent four-state deterministic fixture makes the inconsistency explicit. Let the sole action implement \(T(s)=(s+1)\bmod4\). Suppose training has just observed \(0\to1\). An evaluation segment resets to zero and stops at state three. The cached training observation remains one. The next actual transition is \(3\to0\), while the learner can record it as \(1\to0\); the valid transition from one is \(1\to2\). Thus the recorded transition has probability zero under the stated transition kernel. A real environment that forbids stepping after an episode ends can instead raise an error. The exact symptom depends on its wrapper, but the uncoordinated mutation remains.\par
This fixture checks the object-sharing mechanism only. It is not a Crafter simulation, and no author callback or training loop was run.\par
\#\# Scope, severity, and repair\par
The issue is Minor at the supported evidentiary level. It is a concrete local training-integrity defect, but the supplied evidence does not establish its frequency, trajectory consequences, or material impact on the historical learning curves. Both baseline and guided training use the wrapper, so the direction of any comparative bias cannot be inferred. A common defect does not make either trajectory stream valid, yet it also does not demonstrate an unfair advantage for one method.\par
Constructing an independent evaluation environment prevents evaluation from mutating training state. That separation restores the state-consistency requirement for this mechanism. A proper regression would also check wrapper/reset behavior and evaluation seeds. Such an integrated test was not performed under the no-author-code-execution constraint. No corrected PPO result is claimed.\par
\# Remaining scientific scope and overall dependence\par
\#\# What the empirical evidence establishes\par
The paper supplies strategy-specific behavior measurements in addition to scalar returns: saplings supplied to cows, achievement counts, and relevant skill milestones. These are useful checks that guidance affects behavior rather than only a reward label. The medium breeding results are qualified by incomplete success across runs, including three successes in eight PPO runs under the reported schedule. This limitation is visible in the tables and should remain part of the conclusion.\par
The observed feasibility proxy is a maximum normalized shaped score, while formal feasibility is probability mass on implementable policies. The observed value proxy is an achieved maximum reward, while formal value compares optimal returns within strategy supports. Their plots compare heuristic proxies, not calibrated estimators of the formal quantities. Sections B.2.2–B.2.4 acknowledge heuristic estimates and future feedback updates. No theorem about their calibration is established.\par
The N=3 EDE and N=8 PPO summaries and displayed confidence bands are descriptive evidence under the authors’ aggregation choices. Without seed-level data, this review cannot reconstruct the confidence-interval computation, check seed independence, or verify every statement of statistical significance. These unresolved checks are not promoted to supported errors.\par
\#\# Interaction budgets and reproducibility\par
Appendix B.4 states that the reward model uses 4,000 frames from 200 trajectories of pretrained PPO agents and 1,000 expert-demonstration frames. The released collection code refers to 500,000-step and three-million-step PPO checkpoints. Thus the shaping pipeline uses task-specific offline experience in addition to the plotted online training steps. The experiments can fairly measure downstream online sample efficiency given that prior; they do not by themselves establish a smaller total environment-interaction budget from scratch. Neither the checkpoint creation cost nor its amortization across strategies is resolved here.\par
The default released configuration also differs from the experimental settings: it contains a 5,000-step training value, 10,000 annotation samples, GPT-4o, a human-data ratio of one, and a short linear schedule. Command-line overrides are supported, so these defaults do not prove that the paper used the wrong settings. They show why exact run configurations and provenance are necessary for replication. The snapshot refers to missing checkpoint paths and an imported |test\_reward\_function| module not included in the supplied inventory. These are reproduction limits, not evidence that the reported runs never occurred.\par
\#\# Dependence and surviving conclusions\par
Problem 1 affects the claimed automatic realization of formal specificity. Problem 2 affects the operative automatic choice of strategy in the released algorithm. Neither refutes the valid conditional finite-policy analysis. Problem 3 concerns a separate local rollout mechanism and has a more limited demonstrated impact.\par
The experiments may still support the empirical usefulness of their actual strategy-guided reward functions. Determining how they relate to the paper’s selection algorithm requires historical tree stages, selected-node records, exact configurations, and code versions. Determining revised performance after implementation repairs requires new experiments. No later-version comparison or comprehensive erratum search was conducted; no absence of an erratum is asserted. Author adoption and final scientific adjudication remain pending.\par
\# Compact conclusion\par
Item | Judgment | Repair verification | Remaining work\par
P1: entropy contraction | Major | Conditional/uniform-support correction proved locally | Specify actual policy distributions and operational guarantee\par
P2: strategy priority | Major | Path-product arithmetic checked on supplied trees | Identify historical selection stages and rerun if affected\par
P3: shared evaluation state | Minor | Independent mechanism fixture; separation justified | Integrated wrapper and rollout regression\par
Appendix A | No Issue Found | Discrete support and PAC derivations checked | Does not certify implemented learner\par
Historical costs and selection | Needs Evidence | No experimental repair claimed | Run records, seed data, full budget accounting\par
All entries remain pending human adjudication; no author adoption has been requested.\par
\# Independent checks and source references\par
The reviewer-written computes the entropy counterexample, compares the two priorities on five supplied JSON tree variants, checks the shared-state mechanism, records selected published scalar comparisons, and hashes inspected inputs. It was run with the existing Windows Python 3.13.5:\par
{\ttfamily\footnotesize [local archive path omitted]\\
  ...\textbackslash{}verification\textbackslash{}0198\textbackslash{}independent\_review\_checks.py}\par
The process exited zero. Its output reported parent entropy \(0.1119020568909309\), child entropy \(0.6931471805599453\), and easy-tree winners 11 (local) and 2 (root-path). The evidence output is , with SHA-256\par
808f17036845b8e79c146d8e8c7a65d8
a853e0d2005adbfc97fdab14678d7fb1.\par
Original-PDF text extraction used the existing Poppler executable on all 45 pages. The direct Lumen PDF parser had failed; that failure was not represented as successful reading. Source pages 5, 7, 18, 19, 30, and 31 were rendered for independent parent inspection. Reading a rendering path is not itself recorded as visual inspection.\par
The pinned source paths decisive for the findings are , , , , and . The inspected reward pipeline is . Exact hashes and tree results are retained in the evidence JSON.\par
\# Provenance and roles\par
This is newly performed initial review work, not acceptance of an older report. Earlier partial provenance was preserved once under . The native auxiliary agent assigned to papers 0197–0198 performed the present scientific analysis and independent arithmetic. Its exposed inherited identity is GPT-6 Astra Pro; no unexposed backend identifier is asserted. The parent conversation owns independent evidence integration, exact-final-source compilation, and actual report visual checking; completion of those separate checks belongs in provenance. No human review is claimed, and the human-review list is empty.\par
No author module, training script, pickle model, installation, model API, or paid service was executed. The input PDF and source snapshot were preserved unchanged. The required findings, coverage, and provenance JSON files accompany this source under the assigned output root.\par
9 M. Ruiz Luyten, A. Berthon, and M. van der Schaar, Strategic Planning: A Top-Down Approach to Option Generation, Proceedings of the 42nd International Conference on Machine Learning, PMLR 267 (2025), 52258–52302. https://proceedings.mlr.press/v267/ruiz-luyten25a.html. Reviewed version: 45-page proceedings PDF, including appendices.\par
A. Berthon and collaborators, strategist, released repository, commit |0321ac2cfb1b6a3b407299fc8ce41538b7f24004|. https://github.com/antoninbrthn/strategist. Reviewed as pinned local source; historical experiment provenance unverified.\par
\end{document}
